% Nouveau EGA V (cote 25): archive cote 25, site page 13 (PDF page 14).
% typescript. Draft compared: none.
% AI-transcribed by Claude Opus 5.5 (claude-opus-5-5, adaptive thinking, effort high), Message Batches API, batch ega25b (msgbatch_01A7n5YuXubUPNK1zsepi8RG), 2026-11-13; not yet checked by eye.
\scanpage{25/13}
\mspage{\hand{10}}

sont linéairement indépendantes. Or ces différentielles ne sont autres
que les sections de $(\underline{\Omega}^1_{U/k})\otimes_k k(\ )$ suivantes:
\[ df_1, \ldots, df_p , \sum_{\hand{\unsure{i}}} S_i dT_i \ . \]
Comme les $df_i$ sont linéairement indépendants en tout points de U, et
que les $dT_i$ forment une base \inserted{de $\underline{\Omega}^1_{U/S}$} en tout point de U, \xout{\illegible{1}} et à forti-%
ori un système de générateurs, on en conclut immédiatement l'indépendan-%
ce linéaire des quantités écrites en tout point de $U_{k(\hand{\eta})}$ , du moins
lorsque $p\leq n$, i.e. si \inserted{$\underline{E} =$} $\underline{\Omega}^1_{U/k} / \sum \underline{O}_U df_i \neq 0$ : c'est un petit
lemme sur une famille de \unsure{gébérateurs|générateurs} \inserted{$\unsure{a_i}$ $(1 \leq i \leq n)$} d'un Module loc libre $\underline{E}$ non
nul, alors $\sum S_i a_i$ considéré comme section de $\underline{E}\otimes_k k(\hand{\eta})$ ne s'annule
en aucun point. D(autre part, le cas où p=n est trivial, car alors
$Y_{\hand{\eta}} = \emptyset$ .

\uline{Corollaire} 2.13. Soit Z l'ensemble des points de X en lesquels X
n'est pas lisse sur k. Pour que Y soit lisse sur $k(\hand{\eta})$, il f et s
que Z soit fini. En particulier, si X est lisse, il en est de même de
Y .

Résulte de 2.12. et 2.3. Plus généralement, on obtient:

\margin{\unsure{J un corps} \illegible{1} \unsure{à k} \unsure{de car.}}
\uline{\unsure{Théoèème|Théorème}} 2.1\over{3}{4}. Soi\xout{e}nt \unsure{y} un point de $\over{X}{Y}_{\hand{\eta}}$ , x son image dans X. Soit
\inserted[typed]{$P(\underline{A},K)$} \deleted{P} l'une quelconque des propriétés suivantes pour une\xout{\illegible{1}} algèbre
locale noethérienne sur un corps K:

(i) A est géométriquement régulier sur K.

(ii) A est géométriquement $(R_k)$ sur K.

(iii) A est séparable sur K.

(iv) A est \inserted{géométriquement} normal sur K.

\inserted{Alors $P(\underline{O}_{X,x}, k) \Longleftrightarrow P(\mathcal{O}_{Y_\eta,y}, k(\eta))$}
\xout{(x)xxxxxAxxxx} En effet, compte tenu du par.6 , (iii) et (iv)
résultent de (ii) conjugué avec 2.8. (ii). D'autre part, (i)
a déjà été prouvé dans 2.12. , et il reste à ramener (ii) à (i). Pour

\mspage{\hand{13}}
